\subsection{LLM-Driven Interpretation and Natural Language Querying}

The clustering model in Section~4.1.3 and the sequential pattern mining model in Section~4.1.5 generate highly accurate mathematical results, i.e., prototype vectors, support values, and ordered events. None of these can be directly applied to clinical practice. The current subsection describes the two components that bridge the interpretability gap using the large language model: the directed graph agent that converts natural language queries into safe SQL queries, and the prompt for the summarizer that converts raw data into clinical narrative summaries.

\subsubsection*{Problem Statement}

The critical aspect of this automatic process is that the natural language input ($N$) should be mapped to the safe and read-only database query ($Q_{safe}$):

\begin{equation}
    T(N) = Q_{safe}
    \label{eq:translation}
\end{equation}

It is important to note that the query can only include read-only commands in order to avoid any possible negative actions performed on the clinical database by the agent.

A direct use of the LLM is not sufficient for this problem because LLMs are naturally stochastic, so they have a tendency of hallucinating non-existent database tables or carrying out destructive write actions. It is therefore necessary to enclose the stochastic model within a strict, deterministic control loop.

\subsubsection*{The Agent as a Directed Graph}

The LLM Agent is built in a directed graph form with vertices corresponding to the tasks which have to be done and the edges corresponding to the deterministic sequence of actions. Thus, it ensures that the actual control path in the case of generation by the LLM will be generated by the code.

\vspace{1em}

The state passed between the nodes must be part of the below tuple:

\begin{equation}
    S = \langle\, P,\; \Sigma,\; \rho,\; Q,\; R,\; E_m,\; k \,\rangle,
    \label{eq:state}
\end{equation}

where $P$ denotes the user input, $\Sigma$ denotes the schema of the database, $\rho \in \{\textsc{db}, \textsc{medical}, \textsc{other}\}$ denotes the classification label, $Q$ denotes the generated query, $R$ denotes the retrieved data, $E_m$ denotes the error feedback, and $k$ denotes the iteration number.

The vertex set $V = \{\, v_{\mathrm{in}},\; v_{\mathrm{med}},\; v_{\mathrm{gen}},\; v_{\mathrm{val}},\; v_{\mathrm{exec}},\; v_{\mathrm{end}} \,\}$ defines the operational states:

\begin{itemize}
    \item $v_{\mathrm{in}}$ (\emph{classification}): The LLM classifies $P$ into one of three categories  (DB, MEDICAL or OTHER) and rewrites a database question (if it is the case) for clarity
    \item $v_{\mathrm{med}}$ (\emph{medical answer}): Generates a simple medical response without querying the database
    \item $v_{\mathrm{gen}}$ (\emph{query generation}): Drafts a candidate SQL query $Q$ based on the prompt, schema, and any previous error feedback $E_m$.
    \item $v_{\mathrm{val}}$ (\emph{validation}): Evaluates $Q$ against strict criteria: syntactic validity, schema compliance, and read-only safety.
    \item $v_{\mathrm{exec}}$ (\emph{execution}): Runs the validated query against the database to capture the result set $R$. Any runtime error is treated as a validation failure.
    \item $v_{\mathrm{end}}$ (\emph{termination}): Summarizes $R$ into natural language for the user and ends the process.
\end{itemize}

From the entry node $v_{\mathrm{in}}$, the control flow branches deterministically based on the assigned classification label:

\begin{equation}
    v_{\mathrm{in}} \longrightarrow
    \begin{cases}
        v_{\mathrm{gen}} & \text{if } \rho = \textsc{db},\\
        v_{\mathrm{med}} & \text{if } \rho = \textsc{medical},\\
        v_{\mathrm{end}} & \text{if } \rho = \textsc{other}.
    \end{cases}
    \label{eq:branch}
\end{equation}

In database query processing, the bounded retry cycle includes generation, validation, and execution ($v_{\mathrm{gen}} \to v_{\mathrm{val}} \to v_{\mathrm{exec}}$). When validation fails or an error occurs during execution, the particular error message gets placed in $E_m$, the value of $k$ is increased, and the flow proceeds to $v_{\mathrm{gen}}$.

For the bounded retry cycle to terminate, there must be a successful query result or $k$ reaches its maximum value $k_{\max}$:

\begin{equation}
    \text{exit cycle} \iff f_{\mathrm{val}}(Q,\Sigma) = \mathrm{True} \;\;\lor\;\; k \ge k_{\max}
    \label{eq:exit}
\end{equation}

\subsubsection*{LLM Interpretation of Analytical Output}

The second LLM function takes the numeric output of the clustering and pattern mining processes and generates a summary in a form readable by clinicians. To prevent hallucination, the LLM can be used only as a \emph{verbalizer} for deterministic inputs; no invention or computation of any statistics is allowed.

Each cluster is then provided as input to the LLM through a pre-computed feature pack:

\begin{equation}
    F_c = \langle\, \mu_c,\; \pi_c,\; \delta_c,\; \tau_c,\; \omega_c \,\rangle,
    \label{eq:feature_pack}
\end{equation}

where $\mu_c$ aggregates continuous clinical magnitudes (such as age and lab values), $\pi_c$ denotes the prevalences of categorical attributes, $\delta_c$ captures the demographic distributions, $\tau_c$ signifies the distributional mixes over ordinal risk strata, and $\omega_c$ contains the top frequent clinical events. This tuple encodes all necessary statistics about the cluster. The LLM's task is to identify the most discriminative signals within $F_c$ and phrase them into a natural and accurate clinical analysis.

When this exact same principle is applied to the sequences, the output of the LLM process becomes strictly constrained. Use of deterministic inputs ensures high fidelity of the results produced.
